\section{A Token the Router Will Accept}
\label{sec:ch15-a-token-the-router-accepts}
\index{Cosmo Router!authentication block|(}%

Mosaic signs its own tokens with a shared secret. That is not what a deployment
does, and the reason it is what this repository does is that the alternative
needs an identity provider running before anything can be tested, and this
repository is meant to start with \code{docker compose up}. Cosmo takes the
same shortcut: a JWKS entry can name a \code{url} the router fetches keys from,
or carry a \code{secret} directly, and everything downstream of that choice is
identical. Swapping to a real provider is one key in one file.

Getting the secret form to work cost me most of an afternoon, in two pieces,
and both are the kind of failure that leaves nothing useful behind.

\subsection*{The schema wants three keys, not two}

\index{Cosmo Router!jwks oneOf}%
The obvious configuration is a secret and the algorithm it is used with. The
router refuses to start:

\begin{minted}[fontsize=\footnotesize,breaklines]{text}
- at '/authentication/jwt/jwks/0': oneOf failed, none matched
  - at '/authentication/jwt/jwks/0': validation failed
    - at '/authentication/jwt/jwks/0': missing property 'url'
    - at '/authentication/jwt/jwks/0': not failed
  - at '/authentication/jwt/jwks/0': missing property 'header_key_id'
\end{minted}

This is a good error once you know how to read it, and it is worth reading,
because the \code{oneOf} it is complaining about is a genuine fork rather than
a list of optional keys. The \code{url} branch requires \code{url} and forbids
all three of \code{secret}, \code{symmetric\_algorithm} and
\code{header\_key\_id}. The secret branch requires all three of those and
forbids \code{url}, \code{algorithms}, \code{refresh\_interval} and
\code{refresh\_unknown\_kid}. There is no middle position, and half the keys
that look general are only available on one side.

\code{header\_key\_id} is the interesting one. It is the \code{kid} the router
expects to find in the token's own header, so choosing the secret form is also
choosing to mint tokens that carry a key id. Nothing says so.

\subsection*{The template that is not a template}

\index{Cosmo Router!os.ExpandEnv in config.yaml}%
With the third key added the router starts, reports itself healthy, logs
nothing alarming, and rejects every token I minted with \code{unauthorized} and
HTTP 401. No reason, at \code{info} or at \code{debug}.

The signing key in \code{router/config.yaml} was written like this:

\begin{minted}[fontsize=\footnotesize,breaklines]{yaml}
      - secret: "{{.Env.MOSAIC_JWT_SECRET}}"
\end{minted}

which is the templating form Cosmo's documentation uses elsewhere, and which I
had copied without checking whether this file supports it. It does not. The
router expands its configuration with Go's \code{os.ExpandEnv}:

\begin{minted}[fontsize=\footnotesize,breaklines]{go}
// Expand environment variables in the config file
// and unmarshal it into the config struct
configYamlData := os.ExpandEnv(string(configFileBytes))
\end{minted}

\code{os.ExpandEnv} knows \code{\$VAR} and \code{\$\{VAR\}} and nothing else.
So the router took the twenty-six literal characters of that template as its
HMAC key, and every real token failed its signature check. The correct form is
\code{"\$\{MOSAIC\_JWT\_SECRET\}"} and it is a one-character-class difference
from a graph that refuses everybody.

What makes this expensive is that nothing in the failure points at it. The
config schema is satisfied, because a signing secret is a string and that is a
string. The router starts. The 401 says \code{unauthorized}, which is what a
401 says when the token is genuinely bad. There is exactly one line anywhere in
the system that is a symptom, and it is this, at startup:

\begin{minted}[fontsize=\footnotesize,breaklines]{text}
Using a short secret for JWKs may lead to weak security. Consider using a longer secret.
\end{minted}

which is advice about key strength. It fires because that particular literal is
under the thirty-two characters the check uses. Write a longer variable name
and the warning goes away and the graph still refuses everybody.

I confirmed the diagnosis rather than guessing it, and the confirmation is the
kind I would want to see in somebody else's book. I signed a token using that
template's own twenty-six characters as the key, sent it, and the graph
answered.

Chapter~\ref{ch:real-time-in-a-federated-world} found the same shape one
chapter ago, where wgc's documented spelling of a subscription key composes
perfectly and does nothing. Different tool, different file, and the same
failure: a configuration value that is syntactically fine, semantically wrong,
and reported by nobody.

\subsection*{The minting script}

\index{mint-token.mjs@\code{mint-token.mjs}}%
A reader cannot exercise any of this without a token, so the repository ships
one script that makes them, with no dependency:

\begin{minted}[fontsize=\footnotesize,breaklines]{text}
node scripts/mint-token.mjs
node scripts/mint-token.mjs --scope "read:pii"
node scripts/mint-token.mjs --scope "read:pii" --subject alice
node scripts/mint-token.mjs --expires-in -60
\end{minted}

An HS256 token is two base64url segments and an HMAC over them, so the file is
short enough to read in one sitting, which matters more here than reaching for
a library would. It writes the \code{kid} the previous subsection made
mandatory, and it writes \code{scope} as a space-separated string, which is
what the router's \code{scope\_claim} reads by default.
The last of those four commands mints a token that expired a minute ago, and
section~\ref{sec:ch15-four-callers-one-query} needs it.

Three properties of that script are deliberate and are the difference between a
fixture and a hazard. It signs with the same variable everything else reads, so
a mismatch is impossible to arrange by accident. It refuses to run when that
variable is unset, rather than defaulting to something. And it is not an
identity provider: it signs whatever it is asked to sign, including a token
claiming scopes nobody granted, which is exactly what you want when the thing
under test is the refusal.

\subsection*{The service half}

\index{JwtBearer@\code{JwtBearer}!MapInboundClaims}%
All seven services validate the same tokens, and the wiring is one call in the
shared library. Two settings in it are worth naming.

\code{MapInboundClaims = false}. Left at its default, the handler rewrites
short JWT claim names into the long WS-Federation URIs, so \code{sub} arrives
under a \code{schemas.xmlsoap.org} name ending in \code{nameidentifier},
and a rule written against the name in the token does not match anything. The
token is the contract. The names in it are the names a resolver should see.

\code{ClockSkew = TimeSpan.Zero}. The default is five minutes, which is
generous enough to hide an expiry bug in any test that runs in seconds, and
this chapter has a test that runs in seconds.

Whether that call belongs in a shared library at all is the question
chapter~\ref{ch:strangling-the-monolith} settled the line with: do two services
have to agree about this? For the wiring, no. Every service does it identically and none of
them cares what the others do, which makes it platform. For the key, the issuer
and the audience, yes, which is why those are configuration and a constant in
one script rather than a class six projects reference.

\medskip

With all of that in place a token now reaches the router and is believed. What
happens next depends entirely on who is holding it, and there turn out to be
four interesting answers.

\index{Cosmo Router!authentication block|)}%
